\documentclass[conference]{IEEEtran}
\usepackage{cite}
\usepackage{amsmath,amssymb}
\usepackage{graphicx}
\usepackage{booktabs,tabularx,array}
\usepackage{xcolor}
\usepackage[hidelinks]{hyperref}
\urlstyle{same}
\newcolumntype{Y}{>{\raggedright\arraybackslash}X}

\newcommand{\todo}[1]{\textcolor{red}{\textbf{[#1]}}}

\begin{document}

\title{Do Object Classes Need Different Augmentations? Measuring and Exploiting Class-Specific Augmentation Benefit in Multi-Class Detection}
\maketitle

% \begin{abstract}

% \end{abstract}

\begin{IEEEkeywords}
data augmentation, object detection, class-adaptive augmentation, augmentation policy, road damage, bridge inspection
\end{IEEEkeywords}

\section{Introduction}
\label{sec:intro}

Data augmentation is a standard part of training object detectors, yet most pipelines apply each transformation to the whole image and therefore to every object in it. Detection images rarely contain only one kind of object. In road-damage imagery, thin longitudinal and transverse cracks appear beside alligator cracking and potholes: in our frozen split of RDD2022~\cite{arya2024rdd}, 35\% of the training and validation images that contain damage contain more than one damage class. Bridge-inspection images mix colour-defined defects such as rust and efflorescence with shape-defined cracks~\cite{flotzinger2024dacl10k}, and everyday scenes in PASCAL VOC~\cite{everingham2010voc} routinely place people beside vehicles and furniture. Automated road and bridge inspection depends on detecting each of these categories reliably~\cite{shang2024automated,zheng2024deep,vision2025inspection}, so an augmentation recipe should be judged by what it does for each class, not only by its average effect.

There are good reasons to expect that one augmentation does not suit every class. Colour jitter alters the cue that defines rust while leaving the geometry of a crack intact; blur can erase a hairline crack while barely changing a pothole. Class-level evidence points the same way: augmentation that improves average ImageNet accuracy can reduce accuracy on individual classes~\cite{kirichenko2023}, and road-damage studies report class-dependent responses to augmentation~\cite{tsuchiya2019}. In detection, the question becomes a joint decision, because a whole-image edit reaches every object in the image at once, so the choice for one image must account for all the classes it contains. Despite this, whether detection classes differ in \emph{which augmentations help them learn} has seldom been measured directly with paired controls; most of the evidence comes from classification or from aggregate metrics.

Existing augmentation methods choose transformations at three levels of specificity, and none of them answers this question for multi-class images. Shared policies, from AutoAugment, RandAugment and TrivialAugment to their detection counterparts, select one policy for the whole dataset~\cite{cubuk2019auto,cubuk2020rand,muller2021trivial,zoph2020det,chen2021scale}. Class-aware methods condition the policy on the label~\cite{cheung2022adaaug,yoo2023,ahn2023cuda}, and in detection, instance switching adapts per-class sample generation to detection accuracy~\cite{wang2019psis}. Online controllers update the policy during training, but they read different signals. Tolerance-based controllers such as CTAugment and CUDA ask whether the model still recognizes an augmented example~\cite{berthelot2020,ahn2023cuda}. Benefit-based methods such as BDA, DDAS and MADAug ask whether training on augmented data lowers a held-out loss~\cite{lu2023bda,liu2021ddas,hou2023madaug}. Tolerance and benefit can disagree: an edit the detector tolerates easily may teach it little, while a harder edit may teach it more. The benefit-based methods, in turn, score an augmentation with one validation objective pooled over classes, so a gain for one class can mask a loss for another, while the class-conditioned methods are formulated for single-label samples, where one label selects the policy. The methods reviewed in Section~\ref{sec:related} therefore leave open how per-class benefit should be estimated and combined for an image that contains several classes.

This paper first tests whether the premise holds and then whether it can be exploited. The measurement study trains detectors with one augmentation at a time, starting from a shared warm checkpoint, and records each class's change in validation AP50 relative to the base recipe; a two-way analysis over paired seeds tests whether the effect of an edit depends on the class. A second experiment applies an edit only inside one class's boxes, which separates its direct effect on that class from spill-over to other classes through shared parameters. Building on this measurement, we propose \emph{class-utility augmentation}. At regular intervals during training, it estimates for each edit and each class the extra reduction in that class's loss on a held-out reference set that one gradient step on edited rather than unedited images is predicted to produce. This first-order estimate is a gradient alignment, the same quantity that TracIn uses to trace training influence~\cite{pruthi2020tracin}. For each training image, the policy averages these benefits over the classes present, keeps ``no edit'' as a candidate with zero benefit, blocks any edit that is reliably harmful to one of those classes, and samples a single whole-image edit. A shared-utility variant uses identical machinery but pools the benefit across classes, so comparing the two isolates the contribution of class information.

The evaluation covers three domains (road damage, bridge damage, and general objects), two detector families (Faster R-CNN~\cite{ren2017faster} and RetinaNet~\cite{lin2017focal}), and three paired seeds, with intervals from a paired image-level bootstrap. Four hypotheses are declared before the test set is opened: classes differ in their response to at least one edit (H1); class-utility augmentation outperforms its shared-utility counterpart on test macro AP50, the primary claim (H2); it outperforms the better of the base recipe and a uniform-menu policy (H3); and its online benefit estimates rank-agree with the measured per-class effects (H4). The design makes either outcome informative: if classes do not differ, the study provides a controlled, multi-domain negative result together with the protocol that establishes it. The contributions are as follows:

\begin{enumerate}
    \item A controlled measurement of class-specific augmentation effects in multi-class detection, with one-edit-at-a-time and one-class-at-a-time experiments across three domains and two detector families, paired seeds, and image-level bootstrap intervals.
    \item Class-utility augmentation, an online policy that selects one edit per training image from per-class gradient-alignment benefit estimates, with a no-edit option and a per-class harm block, together with a shared-utility counterpart that differs from it only in class information.
    \item A pre-declared evaluation against base, uniform-menu, measured-table and shared-utility policies, including a test of whether the online estimates agree with the measured per-class effects.
\end{enumerate}

\section{Related Work}
\label{sec:related}

\subsection{Augmentation for Damage Detection}

Infrastructure-damage studies have mainly used augmentation to obtain more informative examples from limited annotations. Tsuchiya \emph{et al.}~\cite{tsuchiya2019} compared augmentation methods by their class-level effect on road-damage recognition and selected methods accordingly. Aghayan-Mashhady and Amirkhani~\cite{aghayan2024} studied bounding-box augmentation and adversarially generated samples for road-damage detection, and a Detectron2 system for multi-country road damage set augmentation probabilities for each damage type and country from data exploration~\cite{rdd_detectron2}. Instance-level methods go further by placing damage into new contexts, as in content- and perspective-aware cut-and-paste~\cite{rdd_cutpaste}. These studies show that useful augmentation for damage depends on the damage type, but they fix that dependence before training and evaluate it within one domain. The present study measures the dependence directly and asks whether it holds across road damage, bridge damage and general objects.

\subsection{Shared and Detection-Specific Policies}

Automatic augmentation methods select operations, probabilities and magnitudes for a whole dataset. AutoAugment~\cite{cubuk2019auto} searches this space using validation performance, RandAugment~\cite{cubuk2020rand} reduces it to an operation count and a common magnitude, and TrivialAugment~\cite{muller2021trivial} samples one operation and strength at random, showing that a simple uniform menu is a strong baseline. Detection adds the choice of where to apply an operation: Zoph \emph{et al.}~\cite{zoph2020det} learn policies that include operations restricted to box content, Scale-Aware AutoAug~\cite{chen2021scale} searches image- and box-level policies conditioned on object scale, and InterAug~\cite{thopalli2023interaug} provides a tuning-free detection policy. In all of these, every image of a given scale receives the same distribution of edits regardless of which classes it contains. We therefore include a TrivialAugment-style uniform menu as a baseline and ask whether class information improves on it.

\subsection{Class-Dependent Effects and Class-Aware Policies}

Averaged gains can hide class-level losses. Kirichenko \emph{et al.}~\cite{kirichenko2023} showed that standard augmentation can reduce accuracy for individual ImageNet classes even as average accuracy rises, which complements the class-dependent road-damage results of Tsuchiya \emph{et al.}~\cite{tsuchiya2019}. Several methods condition augmentation on the class. Yoo and Kang~\cite{yoo2023} derive class-specific parameters from an initially trained classifier and decay them on a schedule during retraining. AdaAug~\cite{cheung2022adaaug} learns class- and instance-adaptive policies that are then reused for training. CUDA~\cite{ahn2023cuda} maintains a severity level for each class in long-tailed recognition. In detection, PSIS~\cite{wang2019psis} switches instances of the same class between images and progressively allocates more switching to classes with low detection accuracy. The classification methods condition on one label per sample, so they do not specify which class should govern an image that contains several; PSIS adapts how many samples each class receives rather than which appearance edit helps it learn. Class-utility augmentation conditions the choice of edit on the full set of classes in each image.

\subsection{Online Feedback: Tolerance and Benefit}

Online controllers adapt augmentation to the model during training, and they can be distinguished by the signal they read. Tolerance-based controllers measure whether the current model still recognizes augmented examples. CTAugment, part of ReMixMatch~\cite{berthelot2020}, updates per-operation magnitude statistics from predictions on augmented labeled examples, pooling observations across classes; CUDA~\cite{ahn2023cuda} raises or lowers each class's severity according to how well augmented examples of that class are recognized. TeachAugment~\cite{suzuki2022teach} similarly keeps transformations recognizable to a teacher while making them hard for the target model. Recognition of an augmented example, however, measures what the model can already handle rather than what it gains from learning on that example. Benefit-based methods target the gain itself. BDA~\cite{lu2023bda} rewards an augmentation by the training gain it produces relative to a paired unaugmented reference; DDAS~\cite{liu2021ddas} uses a one-step gradient update to differentiate held-out loss with respect to the augmentation distribution, including for detection; and MADAug~\cite{hou2023madaug} learns an input-dependent policy jointly with the classifier through a bilevel validation objective. Learned online controllers such as RangeAugment~\cite{mehta2022rangeaugment} and AdaAugment~\cite{yang2025adaaugment} adapt magnitudes from task-loss feedback. These benefit-based methods optimize one objective pooled over classes, so a policy that helps the majority classes can be preferred even when it harms a minority class. The first-order benefit estimate used here relates to TracIn~\cite{pruthi2020tracin}, which scores training examples by the alignment between their gradients and the gradient of a test loss; we compute that alignment separately for each class's reference loss and each candidate edit.

\subsection{Positioning}

Table~\ref{tab:policies} summarizes these mechanisms by the signal they use and by whether that signal is class-specific. The ingredients of the proposed method are not new in isolation: class conditioning, online adaptation, held-out benefit signals, and gradient-alignment influence estimates all have precedents. What the reviewed methods leave open is the combination required by multi-class detection: estimating benefit separately for each class and combining those estimates into one decision for an image that contains several classes, with ``no edit'' as an explicit option. Because this combination is useful only if classes actually differ, the paper measures that premise before evaluating the method, and it evaluates the method against a shared-utility counterpart that removes class information and nothing else.

\begin{table*}[t]
\caption{Augmentation Mechanisms Grouped by Feedback Signal and Class Conditioning}
\label{tab:policies}
\centering
\small
\setlength{\tabcolsep}{4pt}
\renewcommand{\arraystretch}{1.16}
\begin{tabularx}{\textwidth}{@{}p{0.14\textwidth}p{0.15\textwidth}p{0.10\textwidth}p{0.10\textwidth}Y@{}}
\toprule
\textbf{Method} & \textbf{Original setting} & \textbf{Signal} & \textbf{Per class} & \textbf{Relation to this study} \\
\midrule
TrivialAugment~\cite{muller2021trivial} & Classification & None & No & Uniform-menu baseline in our comparison. \\
Zoph \emph{et al.}~\cite{zoph2020det} & Detection & Offline search & No & Box-level operations; one policy for the dataset. \\
Scale-Aware AutoAug~\cite{chen2021scale} & Detection & Offline search & No & Conditioned on object scale rather than class. \\
Detectron2 system~\cite{rdd_detectron2} & Road-damage detection & None (set by data exploration) & Yes & Per-class probabilities fixed before training. \\
PSIS~\cite{wang2019psis} & Detection & Validation AP & Yes & Adapts per-class instance switching, not the choice of appearance edit. \\
Yoo and Kang~\cite{yoo2023} & Classification & Offline, then schedule & Yes & One label per sample selects the parameters. \\
AdaAug~\cite{cheung2022adaaug} & Classification & Learned offline & Yes & One label per sample; policy reused after search. \\
CTAugment~\cite{berthelot2020} & Semi-supervised classification & Tolerance & No & Pooled across classes. \\
CUDA~\cite{ahn2023cuda} & Long-tailed classification & Tolerance & Yes & One label per sample selects the severity. \\
BDA~\cite{lu2023bda} & Classification & Benefit & No & Paired training gain, pooled across classes. \\
DDAS~\cite{liu2021ddas} & Classification, detection & Benefit & No & One-step held-out loss, pooled across classes. \\
MADAug~\cite{hou2023madaug} & Classification & Benefit (learned) & Implicit & Input-dependent policy learned from one pooled validation loss. \\
This study & Multi-class detection & Benefit & Yes & Per-class benefit combined over the classes in each image; no-edit option and per-class harm block. \\
\bottomrule
\end{tabularx}
\vspace{3pt}
\begin{minipage}{\textwidth}
\footnotesize
Signal: \emph{tolerance} = whether the current model recognizes augmented examples; \emph{benefit} = whether training on augmented data reduces a held-out loss. CTAugment is the augmentation component of ReMixMatch.
\end{minipage}
\end{table*}

\section{Methodology}
\label{sec:method}

\subsection{Problem Formulation and Overview}
\label{sec:method_problem}

% Role: opening -- formulate detection and the augmentation decision within that task.
We consider road-damage detection from complete images annotated with bounding boxes and damage categories. Let $\mathcal{D}=\{(x_i,Y_i)\}_{i=1}^{N}$ denote the training set, where $Y_i=\{(b_{ij},c_{ij})\}_{j=1}^{n_i}$ contains the boxes and class labels in image $x_i$. The target categories are longitudinal cracks, transverse cracks, alligator cracks, and potholes, following RDD2022~\cite{arya2024rdd}. A detector $f_\theta$ maps each image to predicted damage locations and categories. Within this task, our objective is to select the transformation family and strength for a local damage region according to its class composition and the detector's current response.

% Role: challenge and overview -- connect the spatial constraint to the central method.
Local augmentation requires a consistent decision wherever annotated boxes share pixels. We therefore organize connected overlapping boxes into groups and assign one transformation to each selected group. The framework regulates transformation strength through class-specific detection retention and selects transformation families using class-balanced loss feedback. Separate groups receive policies conditioned on their own classes, while all classes within an overlap group constrain its shared edit. Ground-truth annotations guide this process during training; inference uses the detector alone.

\subsection{Overall Framework}
\label{sec:method_framework}

% Role: method -- define the policy representation and how feedback is shared across images.
The controller maintains an augmentation policy for each class set represented in an overlap group. Denoting this set by $s$, a group containing a single damage category uses $s=\{c\}$, whereas a mixed-class group uses its complete set of categories. For each $s$, the policy stores a strength for each transformation family and a probability of selecting that family. Groups with the same class set share a policy, and their probe measurements are pooled separately for each represented class. With four categories, this representation contains at most 15 class sets and requires no additional policy network.

% Role: method -- describe the complete feedback and training loop.
Figure~\ref{fig:method_framework} summarizes the interaction between the controller and detector. After an initial warm-up, the controller evaluates paired reference and transformed images using a fixed detector state. These probes determine the strongest eligible level of each family from class-specific retention, then assign family-selection probabilities from normalized loss increases. The resulting policies generate local transformations for detector training. Repeating this update at fixed intervals connects the augmentation choices to the detector's changing response throughout optimization.

\begin{figure*}[t]
    \centering
    \includegraphics[width=\textwidth]{figures/methodology_framework.pdf}
    \caption{Class-adaptive augmentation framework. Paired probes measure the response to a local union edit and pool feedback by group class set. Class-specific retention determines eligible family strengths, and class-balanced loss feedback determines family-selection probabilities. The selected transformation is applied once to a sampled group before full-image detector optimization.}
    \label{fig:method_framework}
\end{figure*}

\subsection{Overlap Groups and Local Transformations}
\label{sec:method_groups}
% Role: method -- turn shared pixel support into a deterministic augmentation unit.
Overlap groups define the spatial units to which the controller assigns transformations. We first clip boxes to the image boundary and rasterize them using the same pixel-support convention as the edit masks. A graph then connects boxes whose rasterized regions have positive-area intersection. Its connected components form the groups, with an isolated box forming a group of its own. Grouping is based on the annotated box regions, including the pavement context they contain.

% Role: method and advantage -- apply a consistent edit across the complete group union.
For a selected group $G$, the transformation acts on the union $M_G$ of its box regions. Let $T_{o,m}$ denote operation $o$ at discrete strength $m$. The augmented image is
\begin{equation}
    \widetilde{x}=x+\alpha_G\odot\bigl(T_{o,m}(x)-x\bigr),
    \label{eq:local_edit}
\end{equation}
where $\odot$ denotes pixelwise multiplication and $\alpha_G$ is a blending mask with values in $[0,1]$. The mask is zero outside $M_G$ and tapers inward at the union boundary to blend the transformed content with its surroundings. We compute $T_{o,m}(x)$ from the reference image and composite it once over the union. This construction gives shared pixels a single transformation and makes the result independent of box-processing order.

% Role: method -- define the transformation bank and the meaning of strength.
The transformation bank contains two coordinate-preserving families: photometric transformations, comprising brightness, contrast, and gamma adjustment; and blur/noise transformations, comprising Gaussian blur and additive Gaussian noise. Opposite directions, such as brightening and darkening, are represented as separate operations. Each operation has a fixed mapping from ordered levels $m\in\{1,\ldots,M\}$ to its physical parameters, with level zero denoting identity. The same parameter mappings, noise distribution, intensity-clipping rule, and boundary taper are used during probing and training. Because these operations preserve spatial coordinates, the original bounding boxes remain the detection targets.

\subsection{Paired Detection Feedback}
\label{sec:method_feedback}
% Role: method -- isolate candidate effects through paired measurements at one detector state.
Paired probes measure the detector's response to each candidate transformation under a common model state. At each policy update, we select a panel of target-bearing training images, preferring distinct source-image clusters, and sample one group uniformly within each image. Each source-image cluster comprises an image and any linked exact or perceptual duplicate candidates; it is distinct from the within-image overlap group. Panel selection is independent of detection outcomes, and measurements are pooled by the selected groups' class sets. The same panel evaluates the operation--strength grid, with only the target group transformed. Probes run in evaluation mode without gradient updates, keeping detector parameters, normalization statistics, preprocessing, and post-processing fixed. For noise, a fixed set of seeded draws is reused across strength levels.

% Role: method -- define class-specific retention of reference detections.
Detection retention measures the fraction of reference-detected objects that remain detected after transformation. Predictions are matched one-to-one to annotated boxes using the correct class label and fixed confidence and intersection-over-union thresholds. Let $d_j^0$ indicate whether reference object $j$ is detected, and let $d_j^{o,m}$ denote its transformed detection indicator, averaged over draws for a stochastic operation. For class-$c$ objects $\mathcal{P}_{s,c}$ in the target groups of class set $s$, retention is
\begin{equation}
 R_{c\mid s,o,m}=
 \frac{\sum_{j\in\mathcal{P}_{s,c}}d_j^0d_j^{o,m}}
      {\sum_{j\in\mathcal{P}_{s,c}}d_j^0}.
 \label{eq:retention}
\end{equation}
Conditioning on reference detections separates missed detections introduced by the edit from those already present in the reference view.

% Role: method -- specify evidence support and complementary full-image measurements.
The controller uses a class-set policy only when every represented class has the required minimum numbers of reference-detected objects and contributing source-image clusters within that exact class set. Repeated noise draws do not increase either count. Measurements from other class sets cannot supply missing support. An absent or insufficiently supported class set receives identity, and its policy is recomputed from the next panel rather than carried forward from an earlier update. Per-object detection indicators, matched-object IoUs, and full-image false-positive counts are retained alongside the target-group measurements to characterize responses beyond the edited region.

\subsection{Joint Strength Selection}
\label{sec:method_strength}
% Role: method -- regulate family strength using the joint class and operation constraint.
Strength selection uses a joint retention constraint across the classes sharing a group. Let $\mathcal{O}_f$ denote the operations in family $f$, including their directions, and let $\mathcal{A}_{s,f}$ be the set of tested nonzero levels with sufficient probe support for which
\begin{equation}
 \min_{c\in s}\;\min_{o\in\mathcal{O}_f}
 R_{c\mid s,o,m}\geq\rho,
 \label{eq:joint_gate}
\end{equation}
where $\rho$ is the retention threshold. The controller selects the strongest level in this admissible set:
\begin{equation}
 m^*_{s,f}=\max\bigl(\mathcal{A}_{s,f}\cup\{0\}\bigr).
 \label{eq:selected_strength}
\end{equation}
The zero value marks a family as unavailable when no tested level satisfies the constraint. Each level is evaluated independently; selecting the strongest admissible level does not assume that retention decreases monotonically with strength. Both eligibility and strength are recomputed from the current probe measurements at each policy update.

% Role: advantage -- explain why the minimum is taken before family allocation.
The minimum in Eq.~\ref{eq:joint_gate} gives every represented class and operation a separate constraint on the shared strength. For example, a strong blur level is rejected for crack--pothole groups when crack retention falls below $\rho$, even if pothole retention remains high. A weaker blur/noise level is eligible only when both blur and noise satisfy the threshold for both classes. Applying the constraint before family selection prevents an acceptable average from masking a failed class-level check and permits each family to retain its own admissible strength.

\subsection{Class-Balanced Family Selection}
\label{sec:method_allocation}
% Role: method -- measure augmentation difficulty using losses attributable to the target objects.
Family selection uses an auxiliary object-level loss to measure difficulty independently of thresholded detection success. For Faster R-CNN, we assign positive region proposals to annotated objects using the detector's region-of-interest (ROI) matcher and append one ground-truth ROI explicitly assigned to each object. This additional ROI supplies a measurement even when the proposal network misses an object. Each ROI contributes classification cross-entropy plus class-specific Smooth L1 box-regression loss, with equal weights and a Smooth L1 transition parameter of $1/9$; the regression term is summed over the four box coordinates. We average ROI contributions within each object before aggregating objects.

% Role: method and boundary -- separate the auxiliary observer from ordinary predictions and training loss.
Ground-truth ROIs are used only by this auxiliary observer and never enter the predictions used for retention or evaluation. Both views use the same assignment and loss procedure, although their proposal sets may differ. Only losses attributable to target-group objects guide family selection. Background classification and proposal-network losses remain part of ordinary detector optimization. The auxiliary score therefore measures a specific supervised response to the edit, rather than the detector's complete training objective or a guarantee of downstream benefit.

% Role: method -- normalize net class-level changes and aggregate classes equally.
For each class, the difficulty score is the normalized increase in its mean object loss. Let $\bar\ell^0_{c\mid s}$ and $\bar\ell^{o,m}_{c\mid s}$ denote mean losses over the same class-$c$ objects in the reference and transformed views, respectively, with stochastic draws also averaged. We define
\begin{equation}
 A_{c\mid s,o,m}=\operatorname{clip}_{[0,B]}
 \left(\frac{\bar\ell^{o,m}_{c\mid s}-\bar\ell^0_{c\mid s}}
 {\bar\ell^0_{c\mid s}+\delta}\right),
 \label{eq:class_difficulty}
\end{equation}
where $\delta>0$ stabilizes normalization and $B>0$ caps the score. Averaging before clipping retains both positive and negative object-level changes when measuring the net class response. Reference-loss normalization adjusts for differences in class loss scale. The score for family $f$ at its selected strength is
\begin{equation}
 D_{s,f}=\frac{1}{|s|\,|\mathcal{O}_f|}
 \sum_{c\in s}\sum_{o\in\mathcal{O}_f}
 A_{c\mid s,o,m^*_{s,f}}.
 \label{eq:family_score}
\end{equation}
The two averages give equal weight to each represented class and to each operation within the family, so classes with more boxes do not dominate family allocation.

% Role: method -- combine difficulty-weighted sampling with continued exploration.
The sampling distribution combines difficulty weighting with uniform exploration over eligible families. For nonempty $\mathcal{E}_s=\{f:m^*_{s,f}>0\}$,
\begin{equation}
 p(f\mid s)=(1-\varepsilon)
 \frac{\exp(D_{s,f}/\tau)}
      {\sum_{h\in\mathcal{E}_s}\exp(D_{s,h}/\tau)}
 +\frac{\varepsilon}{|\mathcal{E}_s|},
 \label{eq:family_probability}
\end{equation}
where $\tau>0$ controls the concentration of the difficulty-weighted component and $0<\varepsilon\leq1$ sets the uniform exploration weight. The uniform term gives every eligible family a nonzero probability. After selecting a family, we sample one of its operations uniformly and apply it at $m^*_{s,f}$. Stochastic operations use a fresh draw from the same distribution used in probing. If no family is eligible, the group receives the identity transformation.

\subsection{Training Procedure and Evaluation Design}
\label{sec:method_evaluation}
% Role: method -- integrate the selected group edit with full-image detector optimization.
Local augmentation is integrated into training through a fixed application probability $q$. After warm-up, an image selected for augmentation contributes one uniformly sampled group, whose class-set policy determines the family, operation, and strength. Other groups remain unchanged on that visit and can be selected on subsequent visits. This matches the single-group intervention used in probing. The complete image and all its annotations contribute to the detector's classification, localization, and background losses. Images without target boxes receive no local edit and retain background supervision. Policy updates use forward measurements, while detector updates minimize the full detection loss.

% Role: method -- distinguish adaptation data from model-selection and evaluation data.
Policy adaptation uses only training annotations and training-image probes. Validation selects the detector checkpoint and supplies development comparisons; it does not supply the online retention or difficulty measurements. Test data are excluded from optimization, probing, and checkpoint selection. Separate random-number streams govern image order, augmentation decisions, and probe selection. Saving the controller state with the detector, optimizer, and random-number states preserves the adopted policy across interrupted training.

% Role: evidence and boundary -- connect the implemented comparisons to the research question.
The evaluation contrasts identity augmentation, fixed image-level augmentation, fixed local augmentation, online shared local control, and online class-adaptive local control. Section~\ref{sec:experimental_setup} specifies the common detector protocol and the comparisons' scope. Detection retention constrains the measured response of the current model; neither passing this constraint nor increasing the auxiliary loss establishes semantic preservation or improved generalization.

\section{Experimental Setup}
\label{sec:experimental_setup}

\subsection{Dataset and Partitioning}
\label{sec:setup_data}
% Role: setup -- define the data population, exclusions, and target conventions.
We use complete annotated images from the seven RDD2022 source groups: China Drone, China MotorBike, Czech, India, Japan, Norway, and the United States~\cite{arya2024rdd}. The four foreground labels are D00, D10, D20, and D40. Auditing the 38,385 available images excludes one exact duplicate with conflicting annotations and one image with an ambiguous class label; one zero-area target box is removed while retaining its image and remaining targets. The resulting pool contains 38,383 images and 55,002 target objects. Valid difficult and truncated targets are retained as ordinary ground truth. Images without annotations for the four target classes remain in the dataset as background-supervision examples; this designation does not establish that their pavement is damage-free.

% Role: setup and limitation -- explain the frozen split and its independence boundary.
The study uses a custom mixed-source split of 26,863 training, 5,760 validation, and 5,760 reserved test images, containing 38,530, 8,225, and 8,247 target objects, respectively. Exact matches and screened perceptual-similarity links define source-image components that are assigned wholly to one partition. Existing test assignments are preserved, components with prior training exposure remain in training, and unassigned groups are balanced within each source toward a 70/15/15 allocation. Six validation images linked to components with training exposure are reassigned to training through this rule. The manifests and exclusions are fixed before detector training. This grouping reduces leakage through identified related images but does not establish independence of all roads or capture sequences, and the split does not measure unseen-country transfer.

\subsection{Detector and Optimization}
\label{sec:setup_detector}
% Role: setup -- specify the common detector and image representation.
All conditions use Faster R-CNN with a ResNet-50 feature pyramid backbone, initialized from the same COCO-pretrained weights. The ROI predictor is replaced with a five-class head for the four damage categories and background. All backbone stages are trainable, with frozen batch-normalization statistics. Each run starts afresh from this initialization rather than continuing a previous condition's checkpoint. Images are decoded in stored orientation, converted to RGB, resized with aspect ratio preserved, and centered in a $640\times640$ canvas with padding value 114. Bounding boxes follow the actual horizontal and vertical resize factors; predictions are mapped back to original image coordinates for evaluation.

% Role: setup -- fix optimization and distinguish learning-rate warm-up from policy warm-up.
Training uses 20 epochs, a batch size of four, and stochastic gradient descent with momentum 0.9 and weight decay $5\times10^{-4}$. The base learning rate is 0.0025, reached through a linear warm-up over 1,000 optimizer steps from $0.001$ times that value. It decreases by a factor of ten after epochs 12 and 16. The complete budget is 134,320 optimizer steps and 537,260 training-image presentations per condition. Augmented conditions use an unaugmented first epoch before enabling their policies. Computation uses FP32 with automatic mixed precision and TF32 disabled. The implementation uses PyTorch 2.3.0, TorchVision 0.18.0, and one NVIDIA A10 GPU per run on CCAST.

\subsection{Augmentation Conditions}
\label{sec:setup_conditions}
% Role: comparison -- define the five conditions and the controlled spatial contrast.
Table~\ref{tab:augmentation_conditions} defines the five comparison conditions. Identity applies only the common preprocessing. The fixed image-level and fixed local conditions share the same operation parameters and sampling rule: each family has probability $1/2$, followed by a uniform draw among its operations at level two. Their difference is whether the edit covers the resized image content or one sampled overlap-group union. Across all augmented conditions, only target-bearing images are eligible, with application probability $q=0.5$ after the first epoch; images without target boxes remain unchanged. Thus, the spatial comparison does not also introduce augmentation of background-only images.

\begin{table*}[t]
\caption{Augmentation Conditions Under a Common Detector Training Protocol}
\label{tab:augmentation_conditions}
\centering
\small
\setlength{\tabcolsep}{4pt}
\renewcommand{\arraystretch}{1.12}
\begin{tabularx}{\textwidth}{@{}p{0.21\textwidth}p{0.20\textwidth}p{0.23\textwidth}Y@{}}
\toprule
\textbf{Condition} & \textbf{Spatial scope} & \textbf{Feedback pooling} & \textbf{Strength and family selection} \\
\midrule
Identity & No augmentation & None & Identity only \\
Fixed image-level & Full resized content & None & Level two; uniform families \\
Fixed local & One overlap-group union & None & Level two; uniform families \\
Online shared local & One overlap-group union & All probe groups, by class & Shared family strengths and probabilities \\
Online class-adaptive local & One overlap-group union & Exact group class set, by class & Class-set-specific strengths and probabilities \\
\bottomrule
\end{tabularx}
\end{table*}

% Role: comparison and limitation -- identify what the online contrasts can isolate.
The two online conditions use identical probe panels, group choices, candidate transformations, update schedules, and numerical thresholds. The shared controller pools groups across class sets and applies its retention constraint to all four classes. The class-adaptive controller pools only groups with the same exact class set and constrains each class represented in that set. This contrast changes feedback pooling and policy routing, including the support available to each policy. The fixed-local comparison additionally changes access from one prescribed strength to three candidate strengths, so it cannot isolate online adaptation independently of the candidate range. The fixed recipes are specified before their runs and are not selected by a validation policy search.

\subsection{Transformations and Controller Settings}
\label{sec:setup_controller}
% Role: implementation -- define operation parameters and the shared rendering convention.
The operation bank contains six photometric operations and two blur/noise operations, with the three strength levels in Table~\ref{tab:operation_levels}. Brightness scales RGB values; contrast scales deviations from the full-content mean luminance; gamma applies the exponent to values normalized to $[0,1]$. Noise is independently Gaussian for each pixel and channel. Operations act on stored RGB code values, followed by clipping and rounding to 8-bit intensities. For local edits, clipped boxes are rasterized using floor at the minimum edges and ceiling at the maximum edges. The union mask has opacity 0.5 on its outer eight-neighbor boundary and 1 in its one-pixel-eroded interior; one-pixel-wide regions retain opacity 0.5. Padding remains unchanged.

\begin{table}[t]
\caption{Operation Parameters at the Three Candidate Strengths}
\label{tab:operation_levels}
\centering
\small
\setlength{\tabcolsep}{3pt}
\renewcommand{\arraystretch}{1.12}
\begin{tabularx}{\columnwidth}{@{}Yccc@{}}
\toprule
\textbf{Parameter} & \textbf{Level 1} & \textbf{Level 2} & \textbf{Level 3} \\
\midrule
Brightness down/up & 0.9/1.1 & 0.8/1.2 & 0.7/1.3 \\
Contrast down/up & 0.9/1.1 & 0.8/1.2 & 0.7/1.3 \\
Gamma brighten/darken & 0.9/1.1 & 0.8/1.2 & 0.7/1.3 \\
Gaussian blur radius (px) & 0.4 & 0.8 & 1.2 \\
Gaussian noise std.\ (RGB/255) & 0.01 & 0.02 & 0.03 \\
\bottomrule
\end{tabularx}
\end{table}

% Role: implementation -- specify the probe budget, support, and fixed numerical choices.
Controller updates occur before epochs 2, 4, \ldots, 20, with the selected policies held fixed between updates. Each update uses 256 target-bearing training images in total, not 256 per class set. Images are ranked using keyed pseudorandom draws, prioritizing distinct source-image clusters before filling from unused images, and one group is selected uniformly in each image. Gaussian noise uses two independent draws, each reused across strength levels; the other operations use one view per level. Table~\ref{tab:controller_settings} gives the remaining settings. They are fixed engineering choices for this comparison, rather than validation-optimized estimates. The minimum support prevents decisions based on only one or two detections, but does not supply a statistical precision guarantee. Unsupported or inadmissible class sets receive identity without borrowing a shared policy or resampling another group.

\begin{table}[t]
\caption{Settings Shared by the Two Online Controllers}
\label{tab:controller_settings}
\centering
\small
\setlength{\tabcolsep}{4pt}
\renewcommand{\arraystretch}{1.12}
\begin{tabularx}{\columnwidth}{@{}Yl@{}}
\toprule
\textbf{Setting} & \textbf{Value} \\
\midrule
Unaugmented warm-up & 1 epoch \\
Policy updates before epochs & $2,4,\ldots,20$ \\
Probe panel size per update (total) & 256 images \\
Noise draws per level & 2 \\
Reference detection confidence / IoU & 0.25 / 0.50 \\
Retention threshold $\rho$ & 0.80 \\
Reference support per constrained class & 8 objects, 4 clusters \\
Candidate nonzero strengths $M$ & 3 \\
Loss stabilizer $\delta$ / cap $B$ & 0.1 / 5 \\
Softmax temperature $\tau$ & 0.25 \\
Uniform exploration weight $\varepsilon$ & 0.20 \\
Application probability $q$ & 0.50 \\
\bottomrule
\end{tabularx}
\end{table}

\subsection{Evaluation and Computational Accounting}
\label{sec:setup_evaluation}
% Role: evaluation -- define model selection, metrics, and the held-out boundary.
Evaluation uses unaugmented images and averages precision over the four foreground classes and IoU thresholds $0.50{:}0.05{:}0.95$ (AP50:95). Validation is performed after every epoch, and the checkpoint with the highest validation AP50:95 is selected, retaining the earlier epoch on a tie. The metric implementation is pycocotools 2.0.10; complementary outputs include AP50, AP75, per-class AP, and recall with at most 100 detections. Inference uses a score threshold of 0.001 and class-wise nonmaximum suppression at IoU 0.5. Predicted boxes on images without target annotations are counted separately at confidence 0.25. Validation scores are development estimates subject to checkpoint selection; the reserved test partition is excluded from this comparison.

% Role: limitation -- bound conclusions by replication and the development design.
The comparison uses one initialization and sampling seed, 42, with common data order, optimizer settings, and checkpoint-selection rules. Keyed augmentation and probe draws are separated from the detector's random-number stream, but CUDA execution is not constrained to produce identical trajectories. Consequently, small differences cannot be separated from training variability by this design, and the comparison does not estimate between-seed uncertainty. It also does not independently ablate the retention gate, family-specific strength, or difficulty-based allocation.

% Role: cost -- distinguish equal optimizer budgets from the added cost of online feedback.
The optimizer budget is identical across conditions, while online control adds forward computation. Each probe image requires one reference view and 27 transformed views: seven deterministic operations at three levels and two noise draws at each of three levels. Ten updates over 256 images therefore require 71,680 additional image forwards per online condition, plus the auxiliary ROI-head evaluations. Total run time and probe time are recorded separately, alongside policy support, retention, chosen strengths, family probabilities, and actual edit exposure. Equal augmentation opportunities do not imply equal numbers of modified images, because an online policy can select identity.

\bibliographystyle{IEEEtran}
\bibliography{references}

\end{document}
